% Manuscript scaffold. Populate empirical tables only from data/processed outputs.
\documentclass[11pt,a4paper]{article}
\usepackage[margin=25mm]{geometry}
\usepackage{booktabs}
\usepackage{amsmath}
\usepackage{hyperref}
\title{Explainable Fair Priority Bargaining at a Signalized Intersection}
\author{Research manuscript generated from the reproducible EFPB pipeline}
\date{\today}
\begin{document}
\maketitle

\begin{abstract}
This manuscript reports a controlled study of Explainable Fair Priority Bargaining (EFPB), a transparent allocator of safe pedestrian and vehicle service at an isolated four-arm signalized intersection. EFPB applies safety and emergency constraints, maximum-wait protection, a calibrated near-efficient candidate set, and waiting-debt-weighted Nash bargaining. Each decision emits a typed DecisionTrace from which deterministic factual and replay-validated counterfactual explanations are generated. Traffic outcomes are evaluated in paired-seed SUMO simulations and explanation effects are evaluated in a within-participant decision-only, factual, and factual-plus-counterfactual replay study. Empirical values are intentionally left unfilled until validated runs and approved human-subject data exist.
\end{abstract}

\section{Introduction}
The contribution is not a claim that Nash bargaining at an intersection is new. Prior work already considers pedestrian and vehicle phases as bargaining players. This study instead evaluates whether a starvation-safe, efficiency-bounded allocator can expose faithful, machine-verifiable explanations and whether explanation depth improves objective comprehension and calibrated reliance.

\section{Research questions and hypotheses}
The study tests allocation, robustness, factual comprehension, contrastive understanding, and judgment calibration. H1--H5 are recorded in the research plan and preregistration: factual explanations should improve rationale accuracy; counterfactuals should improve decision-flip accuracy; explanation condition should improve sound-versus-anomalous discrimination; fairness judgments should improve for sound decisions without excusing anomalies; and counterfactuals may incur inspection-time cost.

\section{Method}
\subsection{Intersection and scenarios}
The primary network has four through-only approaches, four crosswalks, and three service phases: $P_{ALL}$, $V_{NS}$, and $V_{EW}$. Demand cells are L/M/H at approximately 25/50/75 percent of calibrated feasible capacity, with targeted emergency, accessibility, sensor-degradation, near-saturation, and temporal-imbalance scenarios. Each manifest is generated once from independent named random streams and reused across controllers.

\subsection{EFPB controller}
For mode $m\in\{p,v\}$, predicted service utility is
\[
u_m(a,t)=\min\left(1,\frac{S_m(a,t)}{\max(1,D_m(t))}\right).
\]
The disagreement point is the safe fallback utility. Waiting debt is
\[
\tilde\omega_m=\pi_m\left[1+\kappa\left(\operatorname{clip}(W_m/T_m^{\max},0,1)\right)^\eta\right],
\qquad \omega_m=\frac{\tilde\omega_m}{\sum_k\tilde\omega_k}.
\]
After safety, emergency, individual-rationality, and maximum-wait filtering, candidates within
\[
L_{\min}+\delta\max(L_{\min},L_{floor})
\]
are retained. The selected plan maximizes
\[
\sum_m\omega_m\log(u_m-d_m+\epsilon).
\]
Vehicle direction uses total-delay max-pressure. All values and predicates are recorded in the DecisionTrace.

\subsection{Baselines}
B0 is fixed-time, B1 is efficiency-only adaptive, B2 is equal-weight/Wang-style bargaining, and B3 is EFPB without explanation. B4 and B5 are the same EFPB decisions rendered with factual and factual-plus-counterfactual interfaces; they reuse traces and do not create separate traffic runs.

\subsection{Explanation fidelity}
Templates consume only trace-bound slots. A claimed binding constraint must be evaluated and within tolerance. A counterfactual is displayed only if replaying the perturbed state through the same controller selects the stated alternative and the perturbation is minimal on the declared grid; otherwise the interface abstains.

\subsection{Evaluation and statistics}
Primary traffic outcomes are total person delay, pedestrian/vehicle 95th-percentile waits, maximum normalized wait, one preregistered fairness endpoint, one efficiency endpoint, and starvation violations. Run-level paired differences are analyzed with controller and scenario effects, bootstrap confidence intervals, and multiplicity correction. HCI outcomes use logistic mixed models for correctness, mixed models for log response time and composites, and cumulative-link mixed models for ordinal ratings. Trust increase alone is not a success criterion.

\section{Results}
\textbf{Not yet available.} Insert only values generated by the validated aggregation and statistical-analysis commands. Surrogate smoke and benchmark outputs are engineering evidence and must not be reported as SUMO findings.

\begin{table}[h]
\centering
\caption{Primary traffic outcomes. Populate from the locked analysis table.}
\begin{tabular}{lrrrr}
\toprule
Controller & Person delay & Pedestrian $W_{95}$ & Vehicle $W_{95}$ & Starvation violations\\
\midrule
B0 & TBD & TBD & TBD & TBD\\
B1 & TBD & TBD & TBD & TBD\\
B2 & TBD & TBD & TBD & TBD\\
EFPB & TBD & TBD & TBD & TBD\\
\bottomrule
\end{tabular}
\end{table}

\section{Discussion}
Interpretation must distinguish predicted efficiency-bound compliance from realized SUMO delay, objective allocative fairness from perceived fairness, and higher trust from calibrated reliance. Null or adverse findings are valid outcomes under the preregistered decision criteria.

\section{Reproducibility}
The repository records resolved YAML, manifest hashes, software versions, run IDs, decision traces, explanation validation, metrics, figures, and analysis outputs. See the root README and the methodology gap report for the execution contract.
\end{document}
